\begin{table}[htbp]
\centering
\small
\setlength{\tabcolsep}{4pt}
\caption{Dokumentovana očekivanja iz perioda standardizacije AES-a i
proverljivi ishodi do 2026. godine. Ocena je opisna i ne predstavlja
numeričko rangiranje.}
\label{tab:retrospektiva}
\begin{tabular}{@{}>{\raggedright\arraybackslash}p{0.15\textwidth} >{\raggedright\arraybackslash}p{0.25\textwidth} >{\raggedright\arraybackslash}p{0.35\textwidth} >{\raggedright\arraybackslash}p{0.15\textwidth}@{}}
\toprule
Tema &
Dokumentovano očekivanje / kriterijum &
Ishod do 2026. &
Retrospektivna ocena \\
\midrule
Bezbednost primitive &
AES proces daje najveću težinu bezbednosti i procenjuje da Rijndael ima
adekvatnu marginu, uz očekivanje daljeg napretka kriptoanalize
\cite{nechvatal2001aes}. &
Related-key i biclique radovi proširili su teorijsku analizu punog AES-a,
ali nisu doveli do praktičnog oporavka ključa u standardnom modelu primene
\cite{biryukov2009aes256,bogdanov2011biclique,nist_ir8319}. &
U skladu sa očekivanjem, uz ogradu. \\

\addlinespace
Softverska efikasnost &
Rijndael je ocenjen kao veoma efikasan na različitim tadašnjim
softverskim platformama i ograničenim uređajima
\cite{nechvatal2001aes}. &
Savremeni softverski AES koristi arhitekturske optimizacije i često
namenske instrukcije; brzina 2026. nije posledica samo generičke
softverske implementacije. &
Delimično u skladu; kontekst promenjen. \\

\addlinespace
Hardverska efikasnost &
Hardverska realizacija i mogućnost paralelnog/pipeline izvršavanja bili
su deo kriterijuma AES procesa \cite{nechvatal2001aes}. &
AES je dobio namensku CPU podršku poput AES-NI i efikasne GPU
implementacije \cite{gueron2012aesni,tezcan2021gpu}. &
U skladu na opštem nivou; mehanizmi su kasniji. \\

\addlinespace
Namenske AES instrukcije &
AES-NI i slične buduće CPU instrukcije nisu bile dokumentovano
predviđanje AES procesa. &
Intel je od Westmere generacije uveo šest AES instrukcija sa
performansnim i implementaciono-bezbednosnim prednostima
\cite{gueron2012aesni}. &
Izvan originalnog obuhvata očekivanja. \\

\addlinespace
Dužine ključa &
Standard mora podržati 128-, 192- i 256-bitne ključeve
\cite{nechvatal2001aes,nist_fips197_2001}. &
Sve tri varijante ostaju deo FIPS~197 i posle ažuriranja iz 2023.
\cite{nist_fips197_2023}. &
U skladu sa kasnijim ishodom. \\

\addlinespace
DES migracija &
AES je razvijen kao naslednik zastarelog DES-a. &
FIPS~46-3 povučen je 19.\ maja 2005, uz preporuku prelaska na AES ili,
za prelazne primene, TDEA \cite{federalregister2005des}. &
U skladu sa kasnijim ishodom. \\

\addlinespace
TDEA kao prelazno rešenje &
TDEA je zadržan kao odobrena prelazna alternativa nakon pojave AES-a. &
Nova primena TDEA zaštite nedozvoljena je posle 31.\ decembra 2023, a
SP~800-67 Rev.~2 povučen je 1.\ januara 2024.
\cite{nist_sp800131ar2,nist_tdea_withdrawal_2023}. &
Duga tranzicija; konačno povučen. \\
\bottomrule
\end{tabular}
\end{table}

\begin{table}[htbp]
\centering
\small
\setlength{\tabcolsep}{4pt}
\ContinuedFloat
\caption[]{Dokumentovana očekivanja i ishodi do 2026. godine (nastavak).}
\begin{tabular}{@{}>{\raggedright\arraybackslash}p{0.15\textwidth} >{\raggedright\arraybackslash}p{0.25\textwidth} >{\raggedright\arraybackslash}p{0.35\textwidth} >{\raggedright\arraybackslash}p{0.15\textwidth}@{}}
\toprule
Tema &
Dokumentovano očekivanje / kriterijum &
Ishod do 2026. &
Retrospektivna ocena \\
\midrule
Veličina bloka &
AES standardizuje blok od 128 bita, nasuprot 64-bitnom DES/TDEA bloku
\cite{nist_fips197_2001}. &
Sweet32 je praktično pokazao ograničenja 64-bitnog bloka pri velikom
obimu podataka; NIST 2024--2026 istražuje i šire 256-bitne blokove za
buduće visokopropusne primene
\cite{bhargavan2016sweet32,nist_sp800197_predraft,nist_wgcm_2026}. &
128 bita je dobro ostarilo, ali nije bez granica. \\

\addlinespace
Implementa\-ciona bezbednost &
AES proces je razmatrao implementacione karakteristike i bočne kanale,
ali u kontekstu tadašnjeg hardvera \cite{nechvatal2001aes}. &
Cache/timing rizici tabelarnih implementacija postali su važan praktičan
problem; AES-NI uklanja lookup tabele i podatkovno zavisno vreme za tu
klasu implementacije \cite{gueron2012aesni}. &
Promenjen implementacioni kontekst. \\

\addlinespace
AEAD i GCM &
GCM nije bio deo originalnog izbora AES primitive; režimi rada
standardizovani su odvojeno. &
SP~800-38D definiše GCM kao AEAD, a TLS~1.3 zahteva implementaciju
AES-128-GCM skupa \cite{nist_sp80038d,rfc8446}. &
Izvan originalnog obuhvata; protokolarni razvoj. \\

\addlinespace
GPU/SIMT izvršavanje &
Savremeni GPU compute model nije bio dokumentovano očekivanje AES procesa. &
AES ima efikasne GPU realizacije, ali kernel propusnost i kompletna
host--device putanja predstavljaju različite metrike
\cite{tezcan2021gpu}. &
Izvan originalnog obuhvata. \\

\addlinespace
Kvantna pretraga &
AES proces je primarno procenjivao klasičnu iscrpnu pretragu i poznatu
klasičnu kriptoanalizu. &
Grover daje kvadratno ubrzanje u idealizovanom modelu, ali konkretni
resursi, serijska dubina i paralelizacija znatno komplikuju praktičnu
procenu \cite{grassl2016grover,jaques2020grover,nist_pqc_faq}. &
Izvan originalnog obuhvata; zahteva ogradu. \\
\bottomrule
\end{tabular}
\end{table}

\FloatBarrier
